In this section, we build on the theory of tangent spaces developed by Cherubini et al.~\cite{diffgeo-article}.
We introduce and compare two coexisting $R$-module structures on the tangent space of an infinitesimally linear type: the pointwise structure and the microlinear structure.
Finally, we conclude the section with a linearization principle: we show that under suitable conditions on the modules, any pointed lifting problem in a commutative square can be solved linearly, provided a non-linear pointed solution exists.

\subsection{Infinitesimal Types and Tangent Spaces}

In synthetic algebraic geometry, tangent spaces are defined in terms of pointed maps out of the first-order neighbourhood of $0$ in the ring $R$.
This subsection introduces the concept of $n$-th order neighbourhoods around a point, which we use to recall the definition of tangent spaces from \cite{diffgeo-article}.

\begin{definition}\label{def:hood-res}
  Let $(X,x)$ be a pointed type.
  The \notion{$n$-th order neighbourhood around $x$} is the subtype $\hood{X,x}{n} \defeq \setbuild{x' : X}{\exists I \substack{\subseteq \\ \mathrm{fg.ideal}} R, I^{n+1} = 0, I = 0 \rightarrow x' = x}$.
\end{definition}

Taking $n$-th order neighbourhoods defines a functor $\hood{\argmhole}{n} : \pType \rightarrow \pType$.
The action on maps $f : (X,x) \rightarrow_\ast (Y,y)$ is given by $\hood{f}{n} \defeq f\mid_{\hood{X,x}{n}}$.

\begin{proposition}\label{prp:hood-fun-welldef}
  $\hood{\argmhole}{n}$ is a well-defined functor.
\end{proposition}

\begin{proof}
  We need to show that for any $x' \in \hood{X,x}{n}$, we have $f(x') \in \hood{Y,y}{n}$.
  By assumption, we merely have a finitely generated ideal $I \subseteq R$ such that $I^{n+1} = 0$ and if $I = 0$ then $x' = x$.
  It follows that $f(x') \in \hood{Y,y}{n}$, as if $I = 0$, then $f(x') = f(x) = y$ by applying $f$ to $x' = x$ and by $f$ being pointed.
\end{proof}

Observe that we have an inclusion $i_x : \hood{X,x}{n} \hookrightarrow X$ natural in $(X,x)$.

\begin{definition}\label{def:hood}
  A type $X$ is an \notion{$n$-th order neighbourhood around $x : X$} if the inclusion $i_x : \hood{X,x}{n} \hookrightarrow X$ is an equivalence of types.
\end{definition}

\begin{example}\label{ex:R[x]/(x^{n+1})}
  Consider the type $X \defeq \Spec(\quot{R[x_1, \dots, x_m]}{\ideal{x_1, \dots, x_m}^{n+1}})$.
  The ideal $\ideal{x_1, \dots, x_m}^{n+1}$ is generated by all monomials of total degree $n+1$.
  A point $f : X$ is an $R$-algebra homomorphism $f : \quot{R[x_1, \dots, x_m]}{\ideal{x_1, \dots, x_m}^{n+1}} \to R$, which is determined by the values $f(x_i) = r_i \in R$ for $i \in \set{1,\dots,m}$.
  The imposed relations require that for any monomial $p(x_1, \dots, x_m)$ of degree $n+1$, $f(p(x_1, \dots, x_m)) = p(r_1, \dots, r_m) = 0$.
  Equivalently, the ideal $I \defeq \ideal{r_1, \dots, r_m} \subseteq R$ satisfies $I^{n+1} = 0$.

  The type $X$ is an $n$-th order neighbourhood around $0$, i.e. the map $i : \hood{X,0}{n} \hookrightarrow X$ is an equivalence.
  To see this, consider $f : X$ given by $(r_1, \dots, r_m)$.
  Let $I \defeq \ideal{r_1, \dots, r_m}$. Then $I^{n+1} = 0$ by construction.
  If $I = 0$, then each $r_i=0$, implying that $f = 0$.
\end{example}

For the case $n=1$ in the preceding example, we introduce a special notation.
We define the \notion{$m$-dimensional first-order infinitesimal disk}, or just the \notion{$m$-disk}, as the type $\D(m) \defeq \Spec(\quot{R[x_1,\dots, x_m]}{\ideal{x_1,\dots,x_m}^2})$.
By duality, we can identify $\D(m)$ with the type $\{ x : R^m \mid x_ix_j = 0 \text{ for all } i,j \}$.
The $m$-disk is a first-order neighbourhood around $0$.
Thus, whenever we consider $\D(m)$ as a pointed type, we will always mean the pointed type $(\D(m), 0)$.

Let $(X,x)$ and $(Y,y)$ be pointed types.
We introduce the notation $(Y,y)^{(X,x)}$ for the type of pointed maps $(X,x) \rightarrow_\ast (Y,y)$.
If the pointing on $X$ is understood from context, then we will simply write $(Y,y)^X$.

\begin{definition}\label{def:tangent}
  Let $(X,x)$ be a pointed type.
  The tangent space of $X$ at $x$ is the type $\tangent{X}{x} \defeq (X,x)^{\D(1)}$.
\end{definition}

The tangent space construction defines a functor $\tangent{}{} : \pType \rightarrow \pType$.
If $f : (X,x) \rightarrow_\ast (Y,y)$ is a pointed map, the derivative of $f$ is the map $\der{f}{x} \defeq f_\ast : \tangent{X}{x} \rightarrow_\ast \tangent{Y}{y}$ induced by post-composition.
Taking tangent spaces is a local construction in the following sense.

\begin{lemma}\label{prp:hood-factorization}
  Let $(X,x)$ and $(Y,y)$ be pointed types.
  If $X$ is an $n$-th order neighbourhood around $x$, then every pointed map $f : (X,x) \rightarrow_\ast (Y,y)$ factors uniquely through $\hood{Y,y}{n}$.
\end{lemma}

\begin{proof}
  This is immediate by functoriality of $\hood{\argmhole}{n}$, \Cref{prp:hood-fun-welldef}, as $f$ factors into $i_y \circ \hood{f}{n} \circ i_x^{-1}$.
  This factorization is unique, since $i_y$ is an inclusion.
\end{proof}

\begin{remark}\label{rem:tangent-factorization}
  Let $(X,x)$ be a pointed type.
  Specializing \Cref{prp:hood-factorization} to $\tangent{X}{x} \defeq (X,x)^{\D(1)}$, we observe that any tangent vector $\nu : \tangent{X}{x}$ has its image entirely in $\hood{X,x}{1}$.
  Thus, $\tangent{X}{x} = \tangent{\hood{X,x}{1}}{x}$.
\end{remark}

% Note that $\D(1) = \Spec(\quot{R[x]}{\ideal{x}^2})$ can be identified with $\sum_{x:R} x^2 = 0$, namely the square-zero elements of $R$.
% We make use of this in the following construction.

\subsection{Module Structures on Tangent Spaces}

Let $(X,x)$ be a pointed type.
If $X$ is an $R$-module and $x = 0$, then we will describe the pointwise $R$-module structure on the tangent space $\tangent{X}{x}$.
In the case that $(X,x)$ is infinitesimally linear (to be defined in \Cref{def:inf-lin}), we can equip $\tangent{X}{x}$ with a second $R$-module structure.
In this subsection, we will compare these module structures when $M$ is an $R$-module such that $(M,0)$ is infinitesimally linear.


We begin by defining the pointwise module structure on the tangent space $\tangent{M}{0}$ of an $R$-module $M$.
Let $\upsilon, \nu$ be tangent vectors in $\tangent{M}{0}$.
The pointwise addition is given by $(\upsilon + \nu)(\varepsilon) = \upsilon (\varepsilon) + \nu (\varepsilon)$ for all $\varepsilon : \D(1)$.
The scalar multiplication is given by $(r \cdot \upsilon)(\varepsilon) = r \cdot \upsilon (\varepsilon)$ for all $r \in R$ and $\varepsilon : \D(1)$.
It is clear that these operations are well-defined and define a module structure on $\tangent{M}{0}$.

\begin{remark}
  Let $M$ and $N$ be $R$-modules.
  Taking the derivative of a pointed map $f : M \to_\ast N$ yields a pointed map $\der{f}{0} : \tangent{M}{0} \to_\ast \tangent{N}{0}$.
  It is only when $f$ is linear that $\der{f}{0}$ is linear with respect to the pointwise module structures.
\end{remark}

We define the \notion{canonical map} $\can : M \rightarrow \tangent{M}{0}$ by $m \mapsto (\varepsilon \mapsto \varepsilon \cdot m)$.
It is clear that $\can$ is $R$-linear with respect to the pointwise module structure and natural with respect to $R$-linear maps.

So far, we have only considered the tangent spaces of types that are already $R$-modules.
However, the tangent space $\tangent{X}{x}$ of a general pointed type $(X,x)$ can also carry a natural $R$-module structure.
The condition under which this occurs, known as \notion{infinitesimal linearity}, is formulated using certain canonical maps between infinitesimal disks.

For any $m,n : \N$, we define the following inclusion:
\begin{align*}
  \inc_{m,n} : \D(m) \vee \D(n) & \rightarrow \D(m+n) \\
  \inl (\varepsilon) & \mapsto (\varepsilon, 0) \\
  \inr (\varepsilon) & \mapsto (0, \varepsilon)
\end{align*}
The diagonal of the $1$-disk is defined as:
\begin{align*}
  D : \D(1) & \rightarrow \D(2) \\
  \varepsilon & \mapsto (\varepsilon,\varepsilon)
\end{align*}

\begin{definition}\label{def:inf-lin}
  Let $(X,x)$ be a pointed type.
  $(X,x)$ is \notion{infinitesimally linear} if for every $m, n : \N$ the induced map
  \begin{align*}
    \inc_{m,n}^* : (X,x)^{\D (m+n)} \rightarrow (X,x)^{\D(m) \vee \D(n)}
  \end{align*}
  is an equivalence.
\end{definition}

Infinitesimal linearity is a local property in the following sense:

\begin{proposition}\label{prp:inf-lin-hood}
  Let $(X,x)$ be a pointed type.
  $X$ is infinitesimally linear at $x$ if and only if $\hood{X,x}{1}$ is infinitesimally linear at $x$.
\end{proposition}

\begin{proof}
  By \Cref{prp:hood-factorization} any map $\D(n) \rightarrow_\ast (X,x)$ factors uniquely through $\hood{X,x}{1}$.
  Thus, we have the desired lifting property for $(X,x)$ if and only if we have it for $\hood{X,x}{1}$.
\end{proof}

Let $(X,x)$ be infinitesimally linear.
The tangent space $\tangent{X}{x}$ carries an $R$-module structure, which we call the microlinear module structure.
Let $\upsilon$ and $\nu$ be tangent vectors in $\tangent{X}{x}$.
The microlinear addition $\upsilon \boxplus \nu$ is defined by the composite $(\inc_{1,1}^\ast)^{-1}[\upsilon,\nu] \circ D$.
In other words, we obtain this map by lifting in the following diagram
\begin{center}
  \begin{tikzcd}
    & \D(1) \vee \D(1) \ar{rd}{[\upsilon,\nu]} \ar{d}{\inc_{1,1}} \\
    \D(1) \ar{r}{D} \ar[bend right]{rr}{\upsilon \boxplus \nu} & \D(2) \ar[dashed]{r}{\exists !} & X
  \end{tikzcd}
\end{center}
The microlinear scalar multiplication is defined by $(r \boxdot \upsilon)(\varepsilon) = \upsilon(r\varepsilon)$ for all $r : R$ and $\varepsilon : \D(1)$.
Indeed, to prove that this defines a module structure, one can adapt the proof of \cite[Theorem 4.2.19]{david-orbifolds}.
We will give a direct proof.

\begin{theorem}\label{thm:tangent-inf-lin}
  Let $(X,x)$ be infinitesimally linear.
  Then the microlinear module structure on $\tangent{X}{x}$ is well-defined. \earlyqed
\end{theorem}

\begin{proof}
  We need to verify the axioms to have an $R$-module structure.
  
  Unitality is clear by definition.
  
  Associativity follows from infinitesimal linearity.
  Let $\phi : \D(m) \to X$ and $\psi : \D(n) \to X$ be two different sums of tangent vectors.
  We let $\phi_i$ and $\psi_j$ denote the restrictions to the $i$-th and $j$-th components of $\phi$ and $\psi$, respectively.
  The sum $\phi \boxplus \psi : \D(m+n) \to X$ is uniquely given by $[\phi_1,\dots,\phi_m,\psi_1,\dots,\psi_n]$.
  Thus $\phi \boxplus \psi = \phi_1 \boxplus \dots \boxplus \phi_m \boxplus \psi_1 \boxplus \dots \boxplus \psi_n$, as desired.

  Distributivity is also clear.
  Let $\upsilon$ and $\nu$ be tangent vectors in $\tangent{X}{x}$, and $r$ be an arbitrary element of $R$.
  Then $((r\boxdot\upsilon)\boxplus(r\boxdot\nu))(\varepsilon)$ is given by $[r \boxdot \upsilon, r \boxdot \nu]$.
  Likewise, $(r \boxdot (\upsilon \boxplus \nu))(\varepsilon)$ evaluates to $(\upsilon \boxplus \nu)(r\varepsilon)$ for each $\varepsilon$ in $\D(1)$.
  Thus the second tangent vector is also given by $[r \boxdot \upsilon, r \boxdot \nu]$, as desired.
\end{proof}

The microlinear module structure is natural in the sense that the derivative of any pointed map between infinitesimally linear types is linear.
The following proposition generalizes \cite[Lemma 2.2.6]{diffgeo-article} to this setting.

\begin{proposition}\label{prp:derivative-inf-lin}
  Let $f : (X,x) \rightarrow_\ast (Y,y)$ be a pointed map between infinitesimally linear types.
  Then $\der{f}{x}$ is $R$-linear with respect to the microlinear module structure.
\end{proposition}

\begin{proof}
  Let $\upsilon$ be a tangent vector in $\tangent{X}{x}$.
  Then we have $f(r\boxdot\upsilon)(\varepsilon) = f(\upsilon)(r\varepsilon) = (r \boxdot f(\upsilon))(\varepsilon)$ for each $\varepsilon$ in $\D(1)$.
  Thus, $\der{f}{x}(r\boxdot\upsilon) = r \boxdot \der{f}{x}(\upsilon)$, as desired.

  Let $\upsilon$ and $\nu$ be tangent vectors in $\tangent{X}{x}$.
  Then both triangles in the diagram
  \begin{center}
    \begin{tikzcd}
      & \D(1) \vee \D(1) \ar{rd}{[f(\upsilon),f(\nu)]} \ar{d}{\inc_{1,1}} \\
      \D(1) \ar{r}{D} \ar[bend right]{rr}{f(\upsilon \boxplus \nu)} & \D(2) \ar[dashed]{r}{\exists !} & X
    \end{tikzcd}
  \end{center}
  commute, because $f \circ [\upsilon,\nu] = [f(\upsilon),f(\nu)]$.
  Thus $\der{f}{x}(\upsilon \boxplus \nu) = \der{f}{x}(\upsilon) \boxplus \der{f}{x}(\nu)$, as desired.
\end{proof}

If $M$ is an infinitesimally linear $R$-module, its tangent space $\tangent{M}{0}$ carries two potentially different $R$-module structures: the pointwise structure and the microlinear structure. We now compare these two structures.

\begin{theorem}\label{thm:tangent-ptw-micro-add}
  Suppose that $M$ is an infinitesimally linear $R$-module.
  Then the pointwise and microlinear addition on $\tangent{M}{0}$ are the same.
\end{theorem}

\begin{proof}
  Let $\upsilon,\nu : \tangent{M}{0}$.
  The pointwise addition $\upsilon + \nu$ factors through $\D(2)$ by $(\varepsilon,\delta) \mapsto \upsilon(\varepsilon) + \nu(\delta)$.
  It is clear that this map is given by $[\upsilon,\nu]$.
  Thus $\upsilon + \nu = \upsilon \boxplus \nu$.
\end{proof}

\begin{remark}\label{rem:ptw-micro-scalar}
  In general, the scalar multiplication given by the pointwise and microlinear module structures differs.
  Let $\nu : \tangent{M}{0}$ be a tangent vector on $M$, $r : R$, and $\varepsilon : \D(1)$.
  Then $(r\boxdot \nu)(\varepsilon) = \nu(r\varepsilon)$ and $(r \cdot \nu)(\varepsilon) = r\nu(\varepsilon)$.
  We can observe that $r\boxdot\nu = r\cdot \nu$ holds if and only if $\nu$ respects scalar multiplication.
\end{remark}

This comparison allows us to show that the canonical map is also linear with respect to the microlinear structure:

\begin{corollary}\label{cor:can-nat-inf-lin}
  Let $(M,0)$ be an infinitesimally linear $R$-module.
  Then the map ${\can : M \rightarrow \tangent{M}{0}}$ is $R$-linear with respect to the microlinear module structure.
\end{corollary}

\begin{proof}
  The map $\can$ is $R$-linear with respect to the pointwise module structure on $\tangent{M}{0}$.

  Firstly, $\can$ respects the microlinear addition.
  To see this, let $\upsilon$ and $\nu$ be tangent vectors in $\tangent{M}{0}$.
  By \Cref{thm:tangent-ptw-micro-add}, we have $\can(\upsilon \boxplus \nu) = \can(\upsilon + \nu) = \can(\upsilon) + \can(\nu) + \can(\upsilon) \boxplus \can(\nu)$. 

  Secondly, $\can$ respects the microlinear scalar multiplication.
  Let $m : M$, $\varepsilon : \mathbb{D}(1)$ and $r : R$.
  Then we have $(r\boxdot\mathrm{can}(m))(\varepsilon) = \mathrm{can}(m)(r\varepsilon) = r\varepsilon m = \varepsilon rm = \mathrm{can}(rm)(\varepsilon)$.
\end{proof}

\begin{remark}
  The pointwise and microlinear module structures on $\tangent{M}{0}$ agree if $\can$ is surjective.
\end{remark}

\subsection{Tangential Modules}

In this subsection, we study tangential modules (\Cref{def:tangential}).
We show that this class of modules is closed under standard operations like finite limits and finite colimits.
We then conclude this subsection by showing that pointed maps between such modules admit a linearization.

\begin{definition}\label{def:tangential}
  An $R$-module $M$ is \notion{tangential} if $\can : M \rightarrow \tangent{M}{0}$ is an equivalence.
\end{definition}

We can generalize $\can$ to a family of canonical maps $\can^k : M^k \to (M,0)^{\D(k)}$ for all $k \geq 1$, defined by
\[
\can^k (m_1, \dots, m_k) (\varepsilon_1, \dots, \varepsilon_k) = \sum_{i=1}^k \varepsilon_i m_i .
\]

\begin{proposition}\label{prp:tangential-implies-gen-can-equiv}
  Let $M$ be an $R$-module. If $M$ is tangential, then $\can^k$ is an equivalence for all $k \geq 1$.
\end{proposition}

\begin{proof}
  We will show this by induction.
  The base case $k=1$ is trivial by assumption.
  Now, assume the statement holds for $k \ge 1$.
  Let $f : \D(k+1) \to M$ be a pointed map.
  Since $M$ is tangential, the map $\delta \mapsto f(0, \dots, 0, \delta)$ is of the form $\delta m_{k+1}$ for a unique $m_{k+1} : M$.
  Consider the pointed map $\hat{f}(\varepsilon_1, \dots, \varepsilon_k, \delta) \defeq f(\varepsilon_1, \dots, \varepsilon_k, \delta) - \delta m_{k+1}$, which satisfies $\hat{f}(0, \dots, 0, \delta) = 0$.
  For any $\delta : \D(1)$, the map $(\varepsilon_1, \dots, \varepsilon_k) \mapsto \hat{f}(\varepsilon_1, \dots, \varepsilon_k, \delta)$ is a pointed map $\D(k) \to M$.
  By the induction hypothesis, this map is uniquely of the form $\sum_{i=1}^k \varepsilon_i g_i(\delta)$ for some functions $g_i : \D(1) \to M$.
  Define $\hat{g_i}(\delta) \defeq g_i(\delta) - m_i$ where $m_i \defeq g_i(0)$.
  Since $M$ is tangential, the maps $\hat{g_i} : \D(1) \to M$ are uniquely of the form $\delta m'_i$ for some $m'_i : M$.
  Thus we have $g_i(\delta) = m_i + \hat{g_i}(\delta) = m_i + \delta m_i'$.
  We can then compute:
  \[ \hat{f}(\varepsilon_1, \dots, \varepsilon_k, \delta) = \sum_{i=1}^k \varepsilon_i (m_i + \delta m'_i) = \sum_{i=1}^k \varepsilon_i m_i + \sum_{i=1}^k \varepsilon_i \delta m'_i. \]
  Since $(\varepsilon_1, \dots, \varepsilon_k, \delta) : \D(k+1)$, we have that $\varepsilon_i \delta = 0$ for all $i$.
  Therefore, the second term vanishes, yielding $\hat{f}(\varepsilon_1, \dots, \varepsilon_k, \delta) = \sum_{i=1}^k \varepsilon_i m_i$, which shows that $f(\varepsilon_1, \dots, \varepsilon_k, \delta) = \sum_{i=1}^{k+1} \varepsilon_i m_i$ for unique $m_1, \dots, m_{k+1} : M$.
  Thus $\can^{k+1}$ is an equivalence.
\end{proof}

\begin{corollary}\label{cor:tangential-inf-lin}
  Let $M$ be a tangential $R$-module. Then $(M,0)$ is infinitesimally linear.
\end{corollary}

\begin{proof}
  We have the following chain of equivalences factoring $\inc_{m,n}$:
  \begin{center}
    \begin{tikzcd}
      M^{\D(m+n)} \ar{r}{\inc_{m,n}} & M^{\D(m)\vee\D(n)} \ar{rd}[right, below]{\simeq} \\
      & & M^{\D(m)}\times M^{\D(n)} \\
      M^{m+n} \ar{uu}{\can^{m+n}}[right]{\simeq} \ar{r}[below]{\simeq} & M^m \times M^n \ar{ru}{\can^m \times \can^n}[right, below]{\simeq}
    \end{tikzcd}
  \end{center}
  The maps $\can^{m+n}$ and $\can^m \times \can^n$ are equivalences by \Cref{prp:tangential-implies-gen-can-equiv}.
\end{proof}

The properties of being tangential and infinitesimally linear are closed under many common categorical constructions.
This will be important to prove that various modules are themselves tangential.
For instance, we will use this in \Cref{ex:power-series-tangential} to show that the polynomial ring $R[x]$ and power series ring $R[[x]]$ are tangential.

\begin{proposition}[{{\cite[Proposition 5.0.1]{david-orbifolds}}}]\label{prp:lim-inf-lin}
  Infinitesimally linear types are closed under limits.
  If $(X,x)$ is infinitesimally linear, then $(X^A, \mathrm{const}_x)$ is infinitesimally linear for every type $A$.
\end{proposition}

\begin{proposition}\label{prp:exp-tangential}
  Let $M$ be a tangential $R$-module.
  Then for any type $U$, the $R$-module $M^U$ is tangential.
\end{proposition}

\begin{proof}
  % One of the mapping spaces is pointed, so there's a tiny thing to check.
  It is straightforward to check that $\tangent{(M^U)}{0} = (\tangent{M}{0})^U$.
  The latter type is equivalent to $M^U$ by assumption,
  and since the module structure on $M^U$ is pointwise, it follows that the
  resulting equivalence $M^U \to \tangent{(M^U)}{0}$ is homotopic to the canonical map.
\end{proof}

To show that tangential and infinitesimally linear $R$-modules are closed under finite limits and finite colimits, we need some lemmas regarding projectivity.

\begin{lemma}\label{lem:disk-projective}
  For any $m : \N$ the types $\D(m)$ are projective among types.
\end{lemma}

\begin{proof}
  Let $f : X \to \D(m)$ be a surjection. The family $\fib_f : \D(m) \to \UU$ is inhabited by assumption. By Zariski local choice, there exists a unimodular vector $f_1, \dots, f_n : \quot{R[x_1,\dots,x_m]}{\langle x_1, \dots, x_m \rangle^2}$ such that each $f_i$ admits a local section $s_i : D(f_i) \to X$. Since $\quot{R[x_1,\dots,x_m]}{\langle x_1, \dots, x_m \rangle^2}$ is a local ring, at least one of the $f_i$ is a unit, say $f_k$. Then $s_k$ is a global section of $f$.
\end{proof}

\begin{lemma}\label{lem:projective-exact}
  Let $(X,x)$ be projective among types.
  The functor $(\argmhole,0)^{(X,x)} : \modcat{R} \rightarrow \modcat{R}$ is exact.
\end{lemma}

\begin{proof}
  Let $M \twoheadrightarrow N$ be an epimorphism of $R$-modules, and suppose we have a pointed map $f : (X,x) \rightarrow_\ast (N,0)$.
  We merely get a lift $\hat{g} : X \rightarrow M$ of $f$ along the epi.
  The map $g = \hat{g} - \hat{g}(0)$ is a pointed lift of $f$, as desired.
\end{proof}

\begin{proposition}\label{prp:tangent-exact}
  The functors $(\argmhole,0)^{\D(m+n)},\ (\argmhole,0)^{\D(m) \vee \D(n)}$ are exact on $\modcat{R}$.
  In particular, $\tangent{}{0}$ is exact on $\modcat{R}$.
\end{proposition}

\begin{proof}
  Immediate by \Cref{lem:disk-projective} and \Cref{lem:projective-exact}.
\end{proof}

\begin{proposition}\label{prp:op-closure-inf-lin-tangential}
  Infinitesimally linear $R$-modules at $0$ are closed under finite colimits and extensions.
  Tangential $R$-modules are closed under finite limits, finite colimits and extensions.
\end{proposition}

\begin{proof}
  This follows from \Cref{prp:tangent-exact}.
  All cases follow the same pattern.
  We will show how to prove this for the case of infinitesimally linear $R$-modules. An analogous proof works for the other cases.
  The key idea is to use the $5$-lemma to deduce that $\inc_{m,n}^\ast$ is an equivalence, as illustrated in the diagrams below.
  \begin{center}
    \begin{tikzcd}
      A^{\D(m+n)} \ar[hook]{d} \ar{r}[above]{\simeq}[below]{\inc_{m,n}^\ast} & A^{\D(m)\vee\D(n)} \ar[hook]{d} \\
      B^{\D(m+n)} \ar{r}[above]{\simeq}[below]{\inc_{m,n}^\ast} \ar[two heads]{d} & B^{\D(m)\vee\D(n)} \ar[two heads]{d} \\
      C^{\D(m+n)} \ar{r}[above]{}[below]{\inc_{m,n}^\ast} & C^{\D(m)\vee\D(n)}
    \end{tikzcd}
    \qquad
    \begin{tikzcd}
      A^{\D(m+n)} \ar[hook]{d} \ar{r}[above]{\simeq}[below]{\inc_{m,n}^\ast} & A^{\D(m)\vee\D(n)} \ar[hook]{d} \\
      B^{\D(m+n)} \ar{r}[above]{}[below]{\inc_{m,n}^\ast} \ar[two heads]{d} & B^{\D(m)\vee\D(n)} \ar[two heads]{d} \\
      C^{\D(m+n)} \ar[]{r}[above]{\simeq}[below]{\inc_{m,n}^\ast} & C^{\D(m)\vee\D(n)}
    \end{tikzcd}
  \end{center}
  \phantom{M} % This is a hack to get the QED square to get on almost the correct line, but is there a better way?
\end{proof}

We will end this section by showing that tangential modules admit a powerful linearization trick.

\begin{theorem}\label{thm:linearize-lift}
  Suppose we have the following square of tangential $R$-modules
  \begin{center}
    \begin{tikzcd}
      A \ar{r}{u} \ar{d}{m} & X \ar{d}{e} \\
      B \ar{r}{v} & Y
    \end{tikzcd}
  \end{center}
  If there merely exists a non-linear but pointed solution $h':(B,0) \rightarrow_\ast (X,0)$ to the lifting problem above, then there merely exists a linear solution as well.
\end{theorem}

\begin{proof}
  Assume we have a non-linear solution $h':B \rightarrow X$.
  Since $\can$ is a natural isomorphism on the given square, it is sufficient to find such a linear lift on the tangent spaces.
  By \Cref{prp:derivative-inf-lin}, $\der{h'}{0}$ is the desired linear map.
\end{proof}

\begin{remark}
  Explicitly, the linear solution that the above proof produces is $h \coloneq \can^{-1} \circ \der{h'}{0} \circ \can$.
\end{remark}
